% =====================================================================
% III. Datos  (~0.75 páginas)
% =====================================================================
\section{Datos}
\label{sec:datos}

\subsection{Sen1Floods11}
Usamos los 446 recortes etiquetados a mano (\emph{dorados}) de
Sen1Floods11~\cite{bonafilia2020sen1floods11}: imágenes Sentinel-1 GRD de $512\times512$
píxeles ($\approx$10\,m) con las polarizaciones VV y VH en dB, y una máscara con tres valores:
agua (1), no agua (0) y sin dato ($-1$). Respetamos las particiones oficiales
(Tabla~\ref{tab:particiones}): train, valid y test provienen de los mismos 10 eventos, mientras
que Bolivia es un evento completo que nunca se ve en entrenamiento.

\begin{table}[t]
\centering
\caption{Particiones usadas. El agua se mide sobre píxeles válidos.}
\label{tab:particiones}
\begin{tabular}{lrl}
\toprule
Partición & Recortes & Uso \\
\midrule
train   & 252 & entrenamiento (9.4\,\% de agua) \\
valid   &  89 & selección de modelo y de variantes \\
test    &  90 & evaluación final en distribución (E3) \\
Bolivia &  15 & evento no visto (E4) \\
Lima    &  64 & evento peruano no visto (E4) \\
\bottomrule
\end{tabular}
\end{table}

Para la ablación de supervisión débil usamos además los recortes \emph{débiles} del conjunto,
cuyas etiquetas se obtuvieron con el umbral de Otsu sobre el propio SAR. Excluimos los de
Bolivia y los que coinciden con un recorte dorado, para evitar fuga de información. Comparada
con la etiqueta a mano en los recortes dorados, esa etiqueta automática solo alcanza un IoU de
0.495, con mucha comisión en zonas áridas (Somalia, Pakistán): es una supervisión
barata pero ruidosa.

\subsection{Evento de Lima (marzo de 2017)}
Construimos un conjunto de evaluación propio con la misma grilla que Sen1Floods11. La imagen es
la escena Sentinel-1B del 17/03/2017 (10:38 UTC, órbita descendente relativa 171, VV+VH),
obtenida ya calibrada y ortorrectificada desde Google Earth Engine. Las pasadas ascendentes
sobre Lima solo traen VV, por lo que no sirven para un modelo de dos canales. La etiqueta es la
capa de máxima extensión potencial de la inundación del producto EMSN038 de Copernicus EMS
(13.58\,km$^2$), delineada con imágenes ópticas WorldView-2 del 18 y 21 de marzo. Solo se
evalúa dentro del área cubierta por esas imágenes (1\,326\,km$^2$). Los cauces de los ríos Rímac
y Chillón que quedan fuera de la inundación se marcan como $-1$, porque son agua permanente que
la etiqueta no considera inundación. Resultan 64 recortes con solo 1.05\,\% de agua. En 15 de
ellos hay píxeles con VV = VH = 0\,dB fuera de la franja del satélite; se marcan como $-1$ al
leer.

\subsection{Preprocesamiento}
Cada banda se recorta a sus percentiles 1 y 99 y se estandariza con media y desviación
calculadas \emph{solo en train} (VV: $-10.42\pm4.08$\,dB; VH: $-17.27\pm4.79$\,dB). Los píxeles
SAR no finitos ($\approx$2.5\,M en train) pasan a 0 en la entrada y a $-1$ en la etiqueta, de
modo que ni se aprenden ni se evalúan. En entrenamiento se toman recortes aleatorios de
$256\times256$ con volteos y rotaciones de 90$^\circ$. Estas transformaciones no interpolan, así
que la máscara conserva sus tres valores. Valid, test y las regiones externas se evalúan con el
recorte completo de $512\times512$, sin aumentos.
